\section{Overview}\label{3.1}

AURA-Bio is implemented as four independent stages, each usable on its own
and each with its own test suite: an \textbf{acquisition} package that
downloads and validates raw source material, a \textbf{preparation}
package that cleans, deduplicates, chunks, embeds, and indexes that
material, a \textbf{runtime} package that answers one question at a time
using the built indexes, and a \textbf{frontend} that exposes the runtime
through a browser interface. This chapter describes each stage in the
order a single question actually travels through them, since that order is
also the order in which the system's design decisions were made and
tested.

\begin{center}
\footnotesize\ttfamily
\begin{tabular}{c}
User Query \\
$\downarrow$ \\
Task-Type Router $\rightarrow$ Subdomain Router \\
$\rightarrow$ Learner-Level Router \\
$\rightarrow$ Teaching-Style Router $\rightarrow$ Context-Level Selector \\
$\downarrow$ \\
Dense Retrieval (BGE) \hspace{0.6cm} BM25 Retrieval (sparse) \\
$\downarrow$ \hspace{2.2cm} $\downarrow$ \\
Reciprocal Rank Fusion \\
$\downarrow$ \\
Cross-Encoder Reranking \\
$\downarrow$ \\
Evidence Selection (per-source cap, prompt budget) \\
$\downarrow$ \\
Local Biology LLM (draft generation) \\
$\downarrow$ \\
Biology Verification Framework \\
(claims, terms, process, arithmetic, biochemistry) \\
$\downarrow$ \\
Structured Critic (development-time) \\
$\downarrow$ \\
Accept or Revise (bounded to two passes) \\
$\downarrow$ \\
Final Answer, Citations, Confidence, Full Trace
\end{tabular}
\end{center}

\section{Corpus Acquisition and Preparation}\label{3.2}

\subsection{Sources}\label{3.2.1}

The retrieval corpus is built exclusively from official, licence-reviewed
sources, so that every chunk can be traced back to a known origin with a
known usage right. Nine sources contribute directly to retrieval, and a
tenth (PubMedQA) is held out purely as evaluation and instruction-tuning
material, deliberately excluded from every retrieval index so that it can
never leak into a benchmark answer.

\begin{table}[h]
\centering
\caption{Retrieval corpus composition by source}
\label{tab:3.1}
\footnotesize
\begin{tabular}{@{}lrrl@{}}
\toprule
\textbf{Source} & \textbf{Shards} & \textbf{Size (GiB)} & \textbf{Dense-embedded} \\
\midrule
PubMed (filtered) & 1{,}262 & 21.799 & no, BM25 only \\
PMC open-access (XML) & 1{,}231 & 2.007 & yes \\
UniProt/Swiss-Prot & 1 & 0.707 & yes \\
Wikipedia (biology categories) & 1{,}693 & 0.297 & yes \\
NLM LitArch & 6{,}941 & 0.147 & yes \\
Gene Ontology & 1 & 0.013 & yes \\
MeSH descriptors & 1 & 0.011 & yes \\
LibreTexts (open textbooks) & 70 & 0.004 & yes \\
Reactome pathway summaries & 1 & 0.002 & yes \\
\midrule
\textbf{Total (retrieval)} & & \textbf{25.0} & \\
PubMedQA (evaluation only) & 1 & 0.001 & excluded \\
\bottomrule
\end{tabular}
\end{table}

Every one of the 3{,}659 acquired files passed a full validation pass checking
compressed-stream integrity, archive safety, XML/JSON well-formedness, and
provider checksums where published with zero non-\texttt{ok} results
before being admitted to the corpus.

\subsection{Cleaning, Deduplication, and Chunking}\label{3.2.2}

Structured sources (PubMed, PMC, UniProt, Gene Ontology, MeSH, Reactome,
Wikipedia) carry their own stable provider identifiers, which are used
directly as deduplication keys. Less structured sources such as open
textbooks are deduplicated using document-, paragraph-, and chunk-level
text fingerprints. Chunk identifiers are computed as a hash of the
document identifier, the chunk's ordinal position, and its text, so that an
identical chunk id can only arise from a genuinely identical passage.

Retrieval chunks target 450 tokens (measured against the embedding
model's own tokenizer) with a 64-token overlap between adjacent chunks, and
never exceed 510 tokens. This limit was set deliberately below the
generic 500-900 token range often recommended for chunking, because the
BGE embedding model used here truncates silently above 512 tokens: a
chunk built to the generic upper bound would have had text discarded from
its embedding without any error being raised.

\subsection{Data Categories}\label{3.2.3}

The recovered corpus is explicitly split into four distinct uses, since
material well-suited to retrieval is not automatically well-suited to
training. A UniProt database record, for instance, is precise and useful
evidence to retrieve, but training a model to imitate its field-value
formatting would make the model worse at writing prose.

\begin{itemize}
    \item \textbf{Knowledge corpus}: the 25.0\,GiB described above, used
        purely for retrieval.
    \item \textbf{Instruction-tuning dataset}: a smaller, curated set of
        question-answer pairs and explanations, described fully in
        Section~\ref{3.9}.
    \item \textbf{Evaluation dataset}: a completely separate 441-question
        benchmark (Section~\ref{3.8}), never used for training.
    \item \textbf{Failure dataset}: records of answers the verification
        or critique layer flagged during development, retained for possible
        future training use once a human-approval step is built.
\end{itemize}

\section{Retrieval Architecture}\label{3.3}

\subsection{Dense Retrieval}\label{3.3.1}

Dense retrieval represents each chunk as a 768-dimensional vector produced
by \texttt{BAAI/bge-base-en-v1.5} \cite{xiao2023bge}, chosen for its
balance of retrieval quality and inference cost. All 10{,}394{,}250
retrieval-eligible chunks were embedded, producing 14.9\,GiB of vectors in
eleven batch jobs run on Kaggle's free GPU tier, each batch recording its
exact input range and a checksum of its output so that a partial run can be
resumed without re-embedding anything already done.

The vectors are indexed with FAISS \cite{johnson2019faiss} using an
\texttt{IVF,PQ} (inverted-file, product-quantised) index rather than a
flat or scalar-quantised one. This was a measured decision, not a default:
the installed FAISS build on the development machine cannot memory-map an
IVF index at all, this is confirmed by testing on an index of only 5{,}000
vectors, so any index that does not fit entirely in RAM cannot be served.
An uncompressed scalar-quantised index over the full corpus would occupy
approximately 7.6\,GiB, which does not comfortably coexist in memory
alongside the query encoder and the reranker on the available hardware. The
product-quantised index occupies roughly 1\,GiB and is built in a single
in-memory pass.

\subsection{Sparse Retrieval}\label{3.3.2}

Sparse retrieval uses BM25 \cite{robertson2009bm25} over a SQLite
FTS5 full-text index, chosen specifically for exact biological
terminology: gene names, enzyme names, abbreviations, where a dense
embedding can blur a distinction a literal keyword match preserves.
Query execution is conjunctive first: all query terms are required to
match, and the query relaxes to a looser \texttt{OR} match only if the
conjunctive form returns too few results. This matters because
\texttt{ORDER BY bm25()} in SQLite must score every matching row before
returning the top results, so an unconstrained \texttt{OR} of common terms
against a large index can force the database to score hundreds of
thousands of rows for a single query, a real performance problem measured
directly during this internship and described further in Chapter~4.

\subsection{Hybrid Fusion and Reranking}\label{3.3.3}

Dense and sparse results are combined through Reciprocal Rank Fusion
across three channels: dense BGE, core BM25, and a separately weighted
PubMed BM25 channel, each item's score being the sum, across channels
it appears in, of one divided by its rank plus a constant. The fused
candidate list is then rescored by a \texttt{BAAI/bge-reranker-base}
cross-encoder, which reads the query and each candidate passage together
rather than comparing precomputed vectors, and is measurably more accurate
at judging relevance at the cost of being too slow to run over every
document in the corpus, which is why it is applied only to the fused
candidates rather than to the whole index.

\subsection{Evidence Selection}\label{3.3.4}

Reranked candidates are trimmed to what the prompt budget and the model's
context window can actually hold, subject to a per-source cap that
prevents one prolific source from crowding out source diversity in the
final answer.

\section{Query Routing}\label{3.4}

Every question passes through five routing decisions before retrieval
begins:

\begin{itemize}
\item \textbf{Task type}: which of thirteen categories (factual recall,
process explanation, comparison, numerical question, misconception
correction, and so on) the question falls into.
\item \textbf{Biology subdomain}: one of eighteen labels, from cell
biology to bioinformatics.
\item \textbf{Learner level}: eight levels from elementary to research.
\item The \textbf{teaching style} to use in the response, one of nine
modes described below.
\item The \textbf{context level}, which selects how much retrieval depth
the question warrants.
\end{itemize}

\subsection{Teaching Styles}\label{3.4.1}

Nine teaching modes are supported, each imposing a distinct answer
structure:

\begin{itemize}
\item Direct answer: a concise response for quick factual questions.
\item Guided explanation: prerequisite, intuition, technical detail,
example, summary, self check.
\item Socratic: guiding questions rather than an immediate answer.
\item Exam mode: definition, mechanism, diagram guidance, common mistakes.
\item Homework support.
\item Revision notes.
\item Misconception repair.
\item Research mode.
\item Teach-back, where the learner is asked to explain the concept back
and their explanation is evaluated.
\end{itemize}

This is what allows the same underlying model to visibly change how it
answers depending on who is asking and why.

\subsection{Context Levels}\label{3.4.2}

Five context levels govern retrieval depth. Level 0 skips retrieval
entirely for simple conversational turns. Level 1 retrieves ten passages
and reranks to five, for straightforward definitions. Level 2 retrieves
thirty and reranks to ten, with evidence compression, for multi-concept
or comparison questions. Level 3 performs hierarchical retrieval, 
ranking whole documents by their best-scoring chunk, then pulling in
neighbouring chunks from the top documents so that a section split across
several chunks reads as continuous context. Level 4 decomposes the query
into sub-questions (for example, splitting ``difference between mitosis
and meiosis'' into separate retrievals for each term), fuses the results,
and applies the Level 3 expansion on top.

\section{Answer Generation}\label{3.5}

The local generator is \texttt{aura-bio:v0}, a persona-adapted build of a
Qwen3.5 4-billion-parameter model, served
through Ollama and running entirely on a single consumer GPU. The prompt
assembled for each question numbers the retrieved evidence so it can be
cited inline as \texttt{[1]}, \texttt{[2]}, and so on, states the
system's grounding and citation requirements, and layers the teaching-mode
structure and learner-level vocabulary guidance on top. Two prompt
behaviours were added specifically in response to measured failures
described in Chapter~4: an explicit instruction that retrieved evidence
should \emph{add to} what the model already knows rather than restrict it,
and a length budget that keeps introductory scaffolding brief so the
model's answer does not run out of length before reaching the substance of
the question.

\section{Verification Framework}\label{3.6}

Every draft answer is checked by a composable set of verifiers before it
is accepted or sent back for revision. Two different kinds of check are
deliberately kept separate, because they carry different authority: a
handful of \emph{deterministic} checkers, which recompute or re-derive a
claim and can therefore be trusted to raise a blocking error, and one
\emph{learned} checker, which can only ever flag something for a human or
a stricter deterministic check to confirm.

\begin{table}[h]
\centering
\caption{Verification checkers implemented}
\label{tab:3.2}
\footnotesize
\begin{tabularx}{\textwidth}{@{}p{3.4cm} X p{2.2cm}@{}}
\toprule
\textbf{Checker} & \textbf{What it does} & \textbf{Authority} \\
\midrule
Claim support (NLI) & Labels each claim \texttt{supported}, \texttt{contradictory}, or \texttt{unsupported} against the retrieved evidence using a DeBERTa-based natural-language-inference model \cite{he2021deberta} & Learned (blocks only on a materially contradicted answer) \\
\addlinespace[3pt]
Terminology \& entities & Misspellings, confusable term pairs, unexpanded abbreviations, invalid binomial names & Deterministic \\
\addlinespace[3pt]
Process order \& location & Compares stated step order against canonical orderings for nine processes (DNA replication, mitosis, glycolysis, etc.) and checks cellular compartment claims & Deterministic \\
\addlinespace[3pt]
Numerical (SymPy) & Recomputes Hardy Weinberg, Mendelian ratios, population growth, Michaelis Menten, dilution, and carrier-probability questions \cite{meurer2017sympy} & Deterministic \\
\addlinespace[3pt]
Sequences (Biopython) & Checks stated codon-to-amino-acid, start/stop-codon, complement, and GC-content claims against the standard genetic code \cite{cock2009biopython} & Deterministic \\
\addlinespace[3pt]
Chemical equations \& units & Balances stated equations by element and flags physiologically impossible unit magnitudes & Deterministic \\
\addlinespace[3pt]
Pathway stoichiometry & Compares stated metabolic yields (net ATP, NADH per cycle turn, etc.) against known constants & Deterministic \\
\addlinespace[3pt]
Entity authority (question-side) & Checks whether a question presupposes an entity absent from a 135{,}941-name lexicon built from the corpus's own naming authorities & Deterministic (informational) \\
\addlinespace[3pt]
Uncertainty & Distinguishes established fact from debated, hypothesis, and speculation, and flags an unsupported absolute statement & Deterministic \\
\bottomrule
\end{tabularx}
\end{table}

The decision rule is simple by design: any deterministic \texttt{error}
blocks the answer and triggers a revision naming the specific defect. A
materially contradicted claim (a meaningful share of the answer's claims,
not merely one among many, disagreeing with the evidence) also blocks.
Everything else is accepted, with weaker signals recorded in the trace for
review rather than used to force a rewrite. This rule was not the first one
implemented, and Chapter~4 describes in detail why the earlier, stricter
version had to be replaced.

\section{External Critique}\label{3.7}

During development, a hosted model reviews generated answers as an
independent second opinion, currently GLM (\texttt{glm-4.7-flash}) via
Z.AI. The critic outputs a structured judgement (supported claims, unsupported claims,
terminology errors, missing concepts, and an accept/revise decision) and
is bound by two safety rules: it can never override a deterministic
verifier error, and it can only ever make a decision \emph{stricter} than
the deterministic layer's, never looser. This keeps the local model as the
actual generator at all times. The external model is a reviewer, never a
substitute author.

\section{Benchmark and Evaluation Design}\label{3.8}

Following the reasoning in Section~\ref{2.5}, the entire evaluation is
scored by machine-checkable answer keys rather than by a second language
model's judgement. A benchmark of 441 questions spans all eighteen
biology subdomains, every supported task type, and learner levels from
high school to research. Each question's answer key states, as
applicable:

\begin{itemize}
\item Required concepts that must appear.
\item Forbidden phrases that must not.
\item A numeric answer with a tolerance, checked independently by the
same SymPy solvers used inside the verifier.
\item A canonical step ordering.
\end{itemize}

Eight questions are deliberately adversarial, presupposing an entity that
does not exist (for example, a fabricated enzyme or a fabricated
taxonomic species). A correct answer to these refuses the premise rather
than inventing detail to satisfy it, which is the benchmark's direct test
of hallucination.

Four systems are compared on identical questions:

\begin{itemize}
\item The \textbf{baseline}: the raw model, no pipeline at all.
\item \textbf{basic\_rag}: hybrid retrieval only, with no reranking,
verification, or critique.
\item \textbf{full\_rag}: the complete pipeline.
\item \textbf{finetuned\_rag}: the full pipeline once a fine-tuned
adapter is available.
\end{itemize}

This four-way structure is what lets a specific claim, for example
``the reranker helped'' or ``verification helped'', be attributed to the
component responsible, rather than only ever reporting one combined
number.

\section{Tools and Technology Stack}\label{3.9}

\begin{table}[h]
\centering
\caption{Core technology stack}
\label{tab:3.3}
\footnotesize
\begin{tabularx}{\textwidth}{@{}lX@{}}
\toprule
\textbf{Layer} & \textbf{Tools} \\
\midrule
Local inference & Ollama, Qwen3.5 4B model \cite{qwen2023technical} \\
Embeddings \& reranking & \texttt{sentence-transformers}, BGE base/reranker \cite{xiao2023bge}, SapBERT \cite{liu2021sapbert} \\
Sparse retrieval & SQLite FTS5 / BM25 \cite{robertson2009bm25} \\
Dense retrieval & FAISS \cite{johnson2019faiss} \\
Deterministic verification & SymPy \cite{meurer2017sympy}, Biopython \cite{cock2009biopython}, \texttt{periodictable}, \texttt{pint} \\
Learned verification & DeBERTa-v3 NLI \cite{he2021deberta} \\
Router training & scikit-learn (logistic regression, SVM, grid search) \\
External critic & Z.AI GLM, Google Gemini (client implemented) \\
Backend & FastAPI \\
Frontend & Vite, React 19, TanStack Query, Tailwind \\
Language \& tooling & Python, PyTorch, Hugging Face Transformers, pytest/unittest \\
\bottomrule
\end{tabularx}
\end{table}

\section{Hardware}\label{3.10}

Development and all evaluation reported in Chapter~4 were carried out on
consumer hardware: a laptop with an NVIDIA GTX 1650 (4\,GiB VRAM) and
approximately 16\,GiB of system RAM, supplemented by Kaggle's free GPU
tier for the large one-off embedding batches described in
Section~\ref{3.3.1}. This constraint shaped several concrete decisions
documented in Chapter~4, the choice of a product-quantised rather than
scalar-quantised dense index, running the cross-encoder reranker and the
NLI verifier on CPU rather than GPU so as not to evict the language model
from the available VRAM, and sizing the fine-tuning plan in Chapter~5
around what is achievable on this hardware rather than around an idealised
budget.
